\begin{exercise}
\noindent\textbf{Exercise 1.1.1 (Boolean closure).}
Show that if $E, F \subset \mathbb{R}^d$ are elementary sets,
then the union $E \cup F$, the intersection $E \cap F$,
and the set theoretic difference
$
E \setminus F := \{x \in E : x \notin F\},
$
and the symmetric difference
$
E \Delta F := (E \setminus F) \cup (F \setminus E)
$
are also elementary. If $x \in \mathbb{R}^d$, show that the translate
$
E + x := \{y + x : y \in E\}
$
is also an elementary set.

\noindent\textbf{Proof.}

Since $E$ and $F$ are elementary sets, they can be written as
finite unions of boxes:
\[
E=\bigcup_{i=1}^{n}B_i,\qquad
F=\bigcup_{j=1}^{m}C_j.
\]

\noindent\textbf{1. Union}

We have
\[
E\cup F
=\left(\bigcup_{i=1}^{n}B_i\right)
\cup\left(\bigcup_{j=1}^{m}C_j\right).
\]
Since this is a finite union of boxes, $E\cup F$ is elementary.

\noindent\textbf{2. Intersection}

We know that
\[
E\cap F
=\bigcup_{i=1}^{n}\bigcup_{j=1}^{m}(B_i\cap C_j).
\]
Let
\[
B_i=I_{i,1}\times\cdots\times I_{i,d},
\quad
C_j=J_{j,1}\times\cdots\times J_{j,d}.
\]
Then
\[
B_i\cap C_j
=\prod_{r=1}^{d}(I_{i,r}\cap J_{j,r}).
\]
Since the intersection of two intervals is either an interval
or empty, $B_i\cap C_j$ is also a box or empty.

Therefore, $E\cap F$ is a finite union of boxes, so it is
elementary.

\noindent\textbf{3. Difference}

First, consider the difference of two boxes $B_i$ and $C_j$.
We can write
\[
B_i\setminus C_j
=\bigcup_{r=1}^{d}
\left(
I_{i,1}\times\cdots\times
(I_{i,r}\setminus J_{j,r})
\times\cdots\times I_{i,d}
\right).
\]
The difference of two intervals can be written as a union
of at most two intervals. Thus, $B_i\setminus C_j$ is a
finite union of boxes and is elementary.

Now,
\[
\begin{aligned}
E\setminus F
&=\left(\bigcup_{i=1}^{n}B_i\right)
\setminus\left(\bigcup_{j=1}^{m}C_j\right)\\
&=\bigcup_{i=1}^{n}\bigcap_{j=1}^{m}
(B_i\setminus C_j).
\end{aligned}
\]
From parts 1 and 2, finite unions and intersections of
elementary sets are elementary. Therefore, $E\setminus F$
is elementary.

\noindent\textbf{4. Symmetric difference}

By definition,
\[
E\triangle F=(E\setminus F)\cup(F\setminus E).
\]
From part 3, both $E\setminus F$ and $F\setminus E$
are elementary. Their union is also elementary by part 1.

Hence, $E\triangle F$ is elementary.


\noindent\textbf{5. Translation}

Let $x=(x_1,\ldots,x_d)\in\mathbb{R}^d$.
For each box
\[
B_i=I_{i,1}\times\cdots\times I_{i,d},
\]
we have
\[
B_i+x
=(I_{i,1}+x_1)\times\cdots\times(I_{i,d}+x_d).
\]
Since translating an interval gives another interval,
$B_i+x$ is still a box.

Also,
\[
E+x=\bigcup_{i=1}^{n}(B_i+x).
\]
This is a finite union of boxes, so $E+x$ is elementary.
\end{exercise}